% Options for packages loaded elsewhere
\PassOptionsToPackage{unicode}{hyperref}
\PassOptionsToPackage{hyphens}{url}
\PassOptionsToPackage{dvipsnames,svgnames,x11names}{xcolor}
%
\documentclass[
]{agujournal2019}

\usepackage{amsmath,amssymb}
\usepackage{iftex}
\ifPDFTeX
  \usepackage[T1]{fontenc}
  \usepackage[utf8]{inputenc}
  \usepackage{textcomp} % provide euro and other symbols
\else % if luatex or xetex
  \usepackage{unicode-math}
  \defaultfontfeatures{Scale=MatchLowercase}
  \defaultfontfeatures[\rmfamily]{Ligatures=TeX,Scale=1}
\fi
\usepackage{lmodern}
\ifPDFTeX\else  
    % xetex/luatex font selection
\fi
% Use upquote if available, for straight quotes in verbatim environments
\IfFileExists{upquote.sty}{\usepackage{upquote}}{}
\IfFileExists{microtype.sty}{% use microtype if available
  \usepackage[]{microtype}
  \UseMicrotypeSet[protrusion]{basicmath} % disable protrusion for tt fonts
}{}
\makeatletter
\@ifundefined{KOMAClassName}{% if non-KOMA class
  \IfFileExists{parskip.sty}{%
    \usepackage{parskip}
  }{% else
    \setlength{\parindent}{0pt}
    \setlength{\parskip}{6pt plus 2pt minus 1pt}}
}{% if KOMA class
  \KOMAoptions{parskip=half}}
\makeatother
\usepackage{xcolor}
\setlength{\emergencystretch}{3em} % prevent overfull lines
\setcounter{secnumdepth}{-\maxdimen} % remove section numbering
% Make \paragraph and \subparagraph free-standing
\ifx\paragraph\undefined\else
  \let\oldparagraph\paragraph
  \renewcommand{\paragraph}[1]{\oldparagraph{#1}\mbox{}}
\fi
\ifx\subparagraph\undefined\else
  \let\oldsubparagraph\subparagraph
  \renewcommand{\subparagraph}[1]{\oldsubparagraph{#1}\mbox{}}
\fi


\providecommand{\tightlist}{%
  \setlength{\itemsep}{0pt}\setlength{\parskip}{0pt}}\usepackage{longtable,booktabs,array}
\usepackage{calc} % for calculating minipage widths
% Correct order of tables after \paragraph or \subparagraph
\usepackage{etoolbox}
\makeatletter
\patchcmd\longtable{\par}{\if@noskipsec\mbox{}\fi\par}{}{}
\makeatother
% Allow footnotes in longtable head/foot
\IfFileExists{footnotehyper.sty}{\usepackage{footnotehyper}}{\usepackage{footnote}}
\makesavenoteenv{longtable}
\usepackage{graphicx}
\makeatletter
\def\maxwidth{\ifdim\Gin@nat@width>\linewidth\linewidth\else\Gin@nat@width\fi}
\def\maxheight{\ifdim\Gin@nat@height>\textheight\textheight\else\Gin@nat@height\fi}
\makeatother
% Scale images if necessary, so that they will not overflow the page
% margins by default, and it is still possible to overwrite the defaults
% using explicit options in \includegraphics[width, height, ...]{}
\setkeys{Gin}{width=\maxwidth,height=\maxheight,keepaspectratio}
% Set default figure placement to htbp
\makeatletter
\def\fps@figure{htbp}
\makeatother

\usepackage{url} %this package should fix any errors with URLs in refs.
\usepackage{lineno}
\usepackage[inline]{trackchanges} %for better track changes. finalnew option will compile document with changes incorporated.
\usepackage{soul}
\linenumbers
\makeatletter
\@ifpackageloaded{caption}{}{\usepackage{caption}}
\AtBeginDocument{%
\ifdefined\contentsname
  \renewcommand*\contentsname{Table of contents}
\else
  \newcommand\contentsname{Table of contents}
\fi
\ifdefined\listfigurename
  \renewcommand*\listfigurename{List of Figures}
\else
  \newcommand\listfigurename{List of Figures}
\fi
\ifdefined\listtablename
  \renewcommand*\listtablename{List of Tables}
\else
  \newcommand\listtablename{List of Tables}
\fi
\ifdefined\figurename
  \renewcommand*\figurename{Figure}
\else
  \newcommand\figurename{Figure}
\fi
\ifdefined\tablename
  \renewcommand*\tablename{Table}
\else
  \newcommand\tablename{Table}
\fi
}
\@ifpackageloaded{float}{}{\usepackage{float}}
\floatstyle{ruled}
\@ifundefined{c@chapter}{\newfloat{codelisting}{h}{lop}}{\newfloat{codelisting}{h}{lop}[chapter]}
\floatname{codelisting}{Listing}
\newcommand*\listoflistings{\listof{codelisting}{List of Listings}}
\makeatother
\makeatletter
\makeatother
\makeatletter
\@ifpackageloaded{caption}{}{\usepackage{caption}}
\@ifpackageloaded{subcaption}{}{\usepackage{subcaption}}
\makeatother
\ifLuaTeX
  \usepackage{selnolig}  % disable illegal ligatures
\fi
\usepackage{bookmark}

\IfFileExists{xurl.sty}{\usepackage{xurl}}{} % add URL line breaks if available
\urlstyle{same} % disable monospaced font for URLs
\hypersetup{
  pdftitle={Lagged Predictions of Next Week Alcohol Use for Personalized Multimodal Interventions},
  pdfauthor={Kendra Wyant; Jiachen (Coco) Yu; Gaylen Fronk; John J. Curtin},
  pdfkeywords={Substance use disorders, Precision mental health},
  colorlinks=true,
  linkcolor={blue},
  filecolor={Maroon},
  citecolor={Blue},
  urlcolor={Blue},
  pdfcreator={LaTeX via pandoc}}

\journalname{Journal of Important Findings}

\draftfalse

\begin{document}
\title{Lagged Predictions of Next Week Alcohol Use for Personalized
Multimodal Interventions}

\authors{Kendra Wyant\affil{1}, Jiachen (Coco) Yu\affil{1}, Gaylen
Fronk\affil{1}, John J. Curtin\affil{1}}
\affiliation{1}{Department of Psychology, University of
Wisconsin-Madison, }
\correspondingauthor{John J. Curtin}{jjcurtin@wisc.edu}


\begin{abstract}
This study found some pretty cool results that have both high impact and
important clinical implications. For example \ldots{}
\end{abstract}




\subsection{Introduction}\label{introduction}

Digital therapeutics can be optimized for jitais

But at times interventions and recommendations may point the user
outside of the app. (Examples). These multimodal interventions are not
available 24/7. They may take time to implement (eg a day, three days, a
week, two weeks).

In a previous study, we demonstrated that it was possible to predict
future alcohol lapses in the next week, day, and hour with high
sensitivity and specificity (area under the receiver operating curves
{[}auROCs{]} of 0.89, 0.90, and 0.93 respectively). The next hour and
next day models are well-positioned to identify changes in lapse risk
factors and recommend just-in-time interventions to address these
immediate risks. The next week model may not have sufficient temporal
specificity to recommend immediate patient action. Its clinical utility
may improve if we shift this coarser window width into the future. This
``time-lagged'' model could provide patients with increased lead time to
implement multi-modal supports that might not be immediately available
to them (e.g., schedule therapy appointment, request support from an AA
sponsor).

In this study, we train models to predict the probability of lapse at
any point during a week window that is shifted 24 hours (one day), 72
hours (three days), 168 hours (one week), or 336 hours (two weeks) into
the future. Our primary goal is to assess how model performance changes
as lag increases.

\subsection{Methods}\label{methods}

\subsubsection{Participants}\label{participants}

We recruited participants in early recovery (1-8 weeks of abstinence)
from moderate to severe alcohol use disorder in Madison, Wisconsin, US
for a three month longitudinal study. We used data from all participants
included in our previous study (N = 151; see
\href{https://osf.io/preprints/psyarxiv/cgsf7}{Wyant et al., in press}
for enrollment, disposition, and demographic information).

\subsubsection{Lapse Labels}\label{lapse-labels}

We predicted lapses occurring in one week windows. Predictions were
updated hourly. Therefore our models provided hour-by-hour predictions
of future probability of an alcohol lapse within a one week window
shifted at various time points (i.e., lags) into the future (no lag and
one day, three day, one week, or two week lagged predictions). The first
prediction for each participant occurs 24 hours + lag time (in hours)
after their study start date. For example, the first prediction for the
one day lagged prediction model would be 48 hours after participants'
start date.

A prediction window was labeled \emph{lapse} if the start date/hour of
any drinking episode fell within that window. A window was labeled
\emph{no lapse} if no alcohol use occurred within that window +/- 6
hours. If no alcohol use occurred within the window but did occur within
6 hours of the start or end of the window, the window was excluded. We
had previously used a more conservative fence around lapse exclusions
(24 hours vs.~6 hours), however, after re-evaluation we deemed this to
be too conservative.

\subsubsection{Feature Engineering}\label{feature-engineering}

Features were calculated using only data collected before the start of
each prediction window. For our no lag models this included all data
prior to the hour of the start of the prediction window. For our lagged
models, the last EMA data used for feature engineering were collected up
to one day, three days, one week, or two weeks prior to the start of the
prediction window.

Following our previous work, features were derived from three data
sources:

\begin{itemize}
\item
  \emph{Demographics}: We created quantitative features for age and
  personal income, and dummy-coded features for sex, race/ethnicity,
  marital status, education, and employment.
\item
  \emph{Day and time of prediction window onset}: We created dummy-coded
  features to indicate time of day and day of week that the prediction
  window began.
\item
  \emph{Previous EMA responses}: We created raw EMA and change features
  for varying scoring epochs before the start of the prediction window
  for all EMA items (i.e., 12, 24, 48, 72, and 168 hours). Raw features
  included min, max, and median scores for each EMA question across all
  EMAs in each epoch for that participant. We calculated change features
  by subtracting the participant's mean score for each EMA question
  (using all EMAs collected before the start of the prediction window)
  from the associated raw feature. We also created raw and change
  features based on the most recent response for each EMA question and
  raw and change rate features from previously reported lapses.
\end{itemize}

\subsubsection{Modeling Decisions}\label{modeling-decisions}

Our primary performance metric for model selection and evaluation will
be auROC.

We will consider four candidate statistical algorithms including elastic
net, XGBoost, regularized discriminant analysis (rda), and single layer
neural networks. This will ensure we are considering algorithms that
differ in terms of bias-variance tradeoff (e.g., parametric vs
non-parametric) to select the model configuration that most closely
represents the true data generating process.

We will use participant-grouped, nested cross-validation for model
training, selection, and evaluation with auROC. We will use 1 repeat of
10-fold cross-validation for the inner loops (i.e., \emph{validation}
sets) and 3 repeats of 10-fold cross-validation for the outer loop
(i.e., \emph{test} sets).

Best model configurations will be selected using median auROC across the
10 validation sets. Final performance evaluation of those best model
configurations used median auROC across the 30 test sets.

\subsection{Results}\label{results}

\subsection{Discussion}\label{discussion}



\end{document}
